\documentclass[11pt]{article}

\usepackage[a4paper,margin=1in]{geometry}
\usepackage[T1]{fontenc}
\usepackage{lmodern}
\usepackage{amsmath,amssymb,mathtools}
\usepackage{hyperref}

\title{Blueprint for an Analytic Schwinger--Keldysh Kernel Pipeline}
\author{calc\_mbcwc project}
\date{2026-05-20}

\newcommand{\dd}{\mathrm{d}}
\newcommand{\ii}{\mathrm{i}}
\newcommand{\Tr}{\operatorname{Tr}}
\newcommand{\sgn}{\operatorname{sgn}}
\newcommand{\diag}{\operatorname{diag}}
\newcommand{\SK}{\mathrm{SK}}
\newcommand{\eff}{\mathrm{eff}}
\newcommand{\thm}{\mathrm{th}}
\newcommand{\R}{\mathrm{R}}
\newcommand{\A}{\mathrm{A}}
\newcommand{\W}{\mathrm{W}}

\begin{document}
\maketitle

\section{Purpose}

The goal is to define the mathematics that the Python pipeline should
implement.  The input is a weakly coupled finite-temperature model, specified by
its quadratic kernel, interaction tensors, statistics, and index structure. The
output is the leading large-time Schwinger--Keldysh ladder kernel and its
largest positive growth exponent
\begin{equation}
  \lambda_L = \max \operatorname{Re}\operatorname{spec} M ,
\end{equation}
where \(M\) is the linearized ladder evolution operator.

The immediate acceptance benchmark is Stanford's weakly coupled Hermitian
matrix \(\Phi^4\) theory \cite{Stanford:2015owe}.  The general pipeline should
reduce to the formulas in Sections~\ref{sec:stanford} and~\ref{sec:numerics}
when that model is selected.

\section{Model Input}

Represent the Lorentzian action as
\begin{equation}
  S[\phi]
  =
  \frac{1}{2}\int_x \phi_I(x) D^{-1}_{IJ}(\partial)\phi_J(x)
  -
  \sum_{\alpha\in \mathcal V}
  \frac{g_\alpha}{n_\alpha!}
  \int_x
  T^\alpha_{I_1\cdots I_{n_\alpha}}
  \phi_{I_1}(x)\cdots \phi_{I_{n_\alpha}}(x).
  \label{eq:model}
\end{equation}
Here \(I,J\) are collective labels: species, matrix indices, internal indices,
and spin/statistics labels. The parser should store:
\begin{itemize}
  \item the quadratic inverse propagator \(D^{-1}_{IJ}(k)\);
  \item interaction arity \(n_\alpha\), coupling \(g_\alpha\), and tensor
  \(T^\alpha\);
  \item index-contraction rules and large-\(N\) power counting;
  \item dispersion branches \(E_a(\mathbf p)\), either exact or from the poles
  of the retarded propagator.
\end{itemize}

\section{Schwinger--Keldysh Doubling}

For each bosonic field introduce contour fields \(\phi_1,\phi_2\), or
equivalently \(r/a\) fields
\begin{equation}
  \phi_1 = \phi_r + \frac{1}{2}\phi_a,
  \qquad
  \phi_2 = \phi_r - \frac{1}{2}\phi_a .
  \label{eq:ra-map}
\end{equation}
The Schwinger--Keldysh action is
\begin{equation}
  S_{\SK}[\phi_r,\phi_a]
  =
  S[\phi_1] - S[\phi_2].
  \label{eq:sk-action}
\end{equation}
For a single monomial \(\phi^n\), the branch-difference identity is
\begin{equation}
  \phi_1^n-\phi_2^n
  =
  \sum_{\substack{s=1\\s\ \mathrm{odd}}}^{n}
  \binom{n}{s}2^{1-s}
  \phi_r^{\,n-s}\phi_a^{\,s}.
  \label{eq:branch-difference}
\end{equation}
Thus every allowed SK interaction vertex contains an odd number of \(a\)-legs.
This is a required implementation check:
\begin{equation}
  S_{\SK}[\phi_r,\phi_a=0]=0.
  \label{eq:unitarity-check}
\end{equation}

With the convention in Eq.~\eqref{eq:model}, differentiating the interaction
part of \(\ii S_{\SK}\) gives the \(r/a\)-basis vertex
\begin{equation}
  V^\alpha_{b_1\cdots b_n;I_1\cdots I_n}
  =
  -\ii g_\alpha\,
  T^\alpha_{I_1\cdots I_n}
  C(b_1,\ldots,b_n),
  \qquad b_i\in\{r,a\},
  \label{eq:sk-vertex}
\end{equation}
where
\begin{equation}
  C(b_1,\ldots,b_n)
  =
  \begin{cases}
    2^{1-s}, & s=\#\{i:b_i=a\}\ \mathrm{odd},\\
    0,       & s\ \mathrm{even}.
  \end{cases}
  \label{eq:vertex-coefficient}
\end{equation}
If \(T^\alpha\) is not fully symmetrized, the implementation must sum over
inequivalent cyclic/orderings specified by the model file.

\section{Thermal Propagators}

For each quasiparticle branch \(a\), define
\begin{equation}
  E_a(\mathbf p) = \sqrt{\mathbf p^2 + m_a^2}
  \quad\hbox{in the scalar benchmark,}
\end{equation}
and use the retarded propagator
\begin{equation}
  G_{\R,a}(k)
  =
  \frac{\ii}{(k^0+\ii0)^2-E_a(\mathbf k)^2-\Pi_{\R,a}(k)} .
  \label{eq:retarded-general}
\end{equation}
Equivalently, in spectral form,
\begin{equation}
  G_{\R,a}(k)
  =
  \ii\int\frac{\dd\omega}{2\pi}
  \frac{\rho_a(\omega,\mathbf k)}{k^0-\omega+\ii0}.
  \label{eq:spectral-retarded}
\end{equation}
The half-thermal-circle Wightman function used for rungs is
\begin{equation}
  \widetilde G_a(k)
  =
  \frac{\rho_a(k)}{2\sinh(\beta k^0/2)}
  \label{eq:half-wightman}
\end{equation}
for bosons. For a free scalar,
\begin{equation}
  \rho(k)
  =
  \frac{\pi}{E_{\mathbf k}}
  \bigl[
    \delta(k^0-E_{\mathbf k})-\delta(k^0+E_{\mathbf k})
  \bigr],
\end{equation}
and therefore
\begin{equation}
  \widetilde G(k)
  =
  \sum_{s=\pm 1}
  \frac{\pi\,\delta(k^0-sE_{\mathbf k})}
       {2E_{\mathbf k}\sinh(\beta E_{\mathbf k}/2)}.
  \label{eq:free-wightman}
\end{equation}

\section{Self-Energy and Width}

The leading large-time ladder problem requires dressed side rails. Near the
quasiparticle pole use
\begin{equation}
  G_{\R,a}(k)
  \simeq
  \frac{\ii Z_a(\mathbf k)}{2E_a(\mathbf k)}
  \left[
    \frac{1}{k^0-E_a(\mathbf k)+\ii\Gamma_a(\mathbf k)}
    -
    \frac{1}{k^0+E_a(\mathbf k)+\ii\Gamma_a(\mathbf k)}
  \right].
  \label{eq:dressed-retarded}
\end{equation}
The width is obtained from the imaginary part of the on-shell retarded
self-energy,
\begin{equation}
  \Gamma_a(\mathbf p)
  =
  -\frac{Z_a(\mathbf p)}
         {2E_a(\mathbf p)}
  \operatorname{Im}
  \Pi_{\R,a}\bigl(E_a(\mathbf p)+\ii0,\mathbf p\bigr).
  \label{eq:width}
\end{equation}

If a field is massless at tree level, the pipeline must first compute the
thermal mass from the static self-energy:
\begin{equation}
  m_{\eff,a}^2
  =
  m_{0,a}^2 + \Pi_{\thm,a}(0,\mathbf 0).
  \label{eq:thermal-mass-general}
\end{equation}
For the Stanford scalar matrix benchmark,
\begin{equation}
  m_{\thm}^2 = \frac{2\lambda}{3\beta^2}.
  \label{eq:stanford-thermal-mass}
\end{equation}

\section{Rung Construction}

The leading growth kernel is built from rung diagrams: two SK vertices, one on
each real-time fold, joined by Wightman lines. Abstractly, for external
on-shell rail species \(A,B\), the rung has the form
\begin{equation}
  R_{AB}(q)
  =
  \sum_{\gamma\in\mathcal R_{AB}}
  \frac{\mathcal S_\gamma}{|\operatorname{Aut}\gamma|}
  \int
  \prod_{\ell\in I_\gamma}
  \frac{\dd^d\ell}{(2\pi)^d}
  \,
  (2\pi)^d
  \delta^{(d)}
  \left(q+\sum_{\ell\in I_\gamma}\sigma_\ell \ell\right)
  \mathcal I_\gamma(q,\{\ell\}).
  \label{eq:general-rung}
\end{equation}
The integrand is
\begin{equation}
  \mathcal I_\gamma
  =
  V^{L,\gamma}_{A;\{\alpha_\ell\}}
  V^{R,\gamma}_{B;\{\beta_\ell\}}
  \mathcal C_\gamma
  \prod_{\ell\in I_\gamma}
  \widetilde G_{\alpha_\ell\beta_\ell}(\ell),
  \label{eq:rung-integrand}
\end{equation}
where \(\mathcal C_\gamma\) is the index/color contraction factor and
\(\mathcal S_\gamma\) includes statistics signs and multiplicities. The graph
enumerator must retain only rungs whose large-\(N\) and weak-coupling order can
contribute to the requested leading limit.

\section{Ladder Equation}

Let \(f_A(\omega,p)\) denote the amputated ladder amplitude with rail species
\(A\). The exact ladder recursion has the schematic form
\begin{equation}
  f_A(\omega,p)
  =
  -G_{\R,A}(p)G_{\R,A}(\omega-p)
  \left[
    h_A(p)
    +
    \sum_B
    \int\frac{\dd^d k}{(2\pi)^d}
    R_{AB}(k-p)f_B(\omega,k)
  \right],
  \label{eq:ladder-recursion}
\end{equation}
where \(h_A\) is the zero-rung source. The large-time growth exponent is
independent of \(h_A\), so the pipeline drops the inhomogeneous term and keeps
the homogeneous equation.

Using the pole-pair replacement
\begin{equation}
  G_{\R,A}(p)G_{\R,A}(\omega-p)
  \longrightarrow
  -
  \frac{\pi\ii Z_A(\mathbf p)^2}{2E_A(\mathbf p)^2}
  \frac{
    \delta(p^0-E_A(\mathbf p))+\delta(p^0+E_A(\mathbf p))
  }{\omega+2\ii\Gamma_A(\mathbf p)},
  \label{eq:pole-pair}
\end{equation}
the amplitude becomes on-shell supported. Write
\begin{equation}
  f_A(\omega,p)
  =
  f_A(\omega,\mathbf p)\,
  \delta\bigl((p^0)^2-E_A(\mathbf p)^2\bigr).
\end{equation}
Then the real-time evolution equation is
\begin{equation}
  \partial_t f_A(t,\mathbf p)
  =
  \sum_B\int
  \frac{\dd^{d-1}\mathbf k}{(2\pi)^{d-1}}
  m_{AB}(\mathbf k,\mathbf p) f_B(t,\mathbf k)
  -
  2\Gamma_A(\mathbf p)f_A(t,\mathbf p).
  \label{eq:general-evolution}
\end{equation}
The on-shell gain kernel is
\begin{equation}
  m_{AB}(\mathbf k,\mathbf p)
  =
  \sum_{\eta=\pm}
  \frac{
    Z_A(\mathbf p)Z_B(\mathbf k)
    R_{AB}\bigl(k_\eta-p\bigr)
  }
  {4E_A(\mathbf p)E_B(\mathbf k)},
  \qquad
  k_\eta-p
  =
  \bigl(E_B(\mathbf k)+\eta E_A(\mathbf p),\mathbf k-\mathbf p\bigr),
  \label{eq:gain-kernel}
\end{equation}
with the precise sign sum fixed by the pole assignments and statistics.

The operator implemented numerically is
\begin{equation}
  (Mf)_A(\mathbf p)
  =
  \sum_B\int
  \frac{\dd^{d-1}\mathbf k}{(2\pi)^{d-1}}
  m_{AB}(\mathbf k,\mathbf p) f_B(\mathbf k)
  -
  2\Gamma_A(\mathbf p)f_A(\mathbf p).
  \label{eq:M-general}
\end{equation}
Finally,
\begin{equation}
  \lambda_L = \max\{\operatorname{Re}\mu:\mu\in\operatorname{spec}M\}.
  \label{eq:lambda-general}
\end{equation}

\section{Stanford \texorpdfstring{\(\Phi^4\)}{Phi4} Benchmark}
\label{sec:stanford}

The benchmark theory is a Hermitian matrix scalar in four spacetime dimensions,
with
\begin{equation}
  \mathcal L
  =
  \frac{1}{2}
  \Tr\left(
    \dot\Phi^2-(\nabla\Phi)^2-m^2\Phi^2-2g^2\Phi^4
  \right),
  \qquad
  \lambda = g^2N .
  \label{eq:stanford-lagrangian}
\end{equation}
For this theory, the leading rung is
\begin{equation}
  R(q)
  =
  48\lambda^2
  \int\frac{\dd^4\ell}{(2\pi)^4}
  \widetilde G(q/2+\ell)\widetilde G(q/2-\ell).
  \label{eq:stanford-rung-integral}
\end{equation}
The corresponding on-shell kernel is
\begin{equation}
  m(\mathbf k,\mathbf p)
  =
  \frac{R(k_+)+R(k_-)}{4E_{\mathbf k}E_{\mathbf p}},
  \qquad
  k_\pm =
  \bigl(E_{\mathbf k}\pm E_{\mathbf p},\mathbf k-\mathbf p\bigr).
  \label{eq:stanford-mkp}
\end{equation}
The width can be eliminated in favor of the same rung:
\begin{equation}
  \Gamma_{\mathbf p}
  =
  \frac{1}{6}
  \int\frac{\dd^3\mathbf k}{(2\pi)^3}
  \frac{\sinh(\beta E_{\mathbf p}/2)}
       {\sinh(\beta E_{\mathbf k}/2)}
  m(\mathbf k,\mathbf p).
  \label{eq:stanford-width-rung}
\end{equation}
Therefore the eigenvalue equation is
\begin{equation}
  \lambda_L f(\mathbf p)
  =
  \int\frac{\dd^3\mathbf k}{(2\pi)^3}
  m(\mathbf k,\mathbf p)
  \left[
    f(\mathbf k)
    -
    \frac{\sinh(\beta E_{\mathbf p}/2)}
         {3\sinh(\beta E_{\mathbf k}/2)}
    f(\mathbf p)
  \right].
  \label{eq:stanford-3d}
\end{equation}

The closed form rung used for validation is
\begin{equation}
  R(q)
  =
  \frac{6\lambda^2}
       {\pi\beta |\mathbf q|\sinh(|q^0|\beta/2)}
  \left[
    \Theta(-q^2-4m^2)
    \log\frac{\sinh x_+}{\sinh x_-}
    +
    \Theta(q^2)
    \log\frac{1-e^{-2x_+}}{1-e^{2x_-}}
  \right],
  \label{eq:stanford-rung-closed}
\end{equation}
where
\begin{equation}
  q^2 = -(q^0)^2+|\mathbf q|^2,
  \qquad
  x_\pm =
  \frac{\beta}{4}
  \left[
    |q^0|
    \pm
    |\mathbf q|
    \sqrt{1+\frac{4m^2}{|\mathbf q|^2-(q^0)^2}}
  \right].
  \label{eq:stanford-xpm}
\end{equation}

\section{One-Dimensional Reduction and Discretization}
\label{sec:numerics}

For rotationally invariant scalar models, set \(P=|\mathbf p|\),
\(K=|\mathbf k|\), \(E_P=\sqrt{P^2+m^2}\). The angular integration uses
\begin{equation}
  \int\dd^3\mathbf k
  =
  2\pi\int_0^\infty K^2\,\dd K
  \int_{|K-P|}^{K+P}\frac{y\,\dd y}{KP},
  \qquad y=|\mathbf k-\mathbf p|.
  \label{eq:angular-measure}
\end{equation}
The Stanford equation becomes
\begin{equation}
  \lambda_L f(P)
  =
  \int_0^\infty \dd K\, m_1(P,K)
  \left[
    f(K)
    -
    \frac{\sinh(\beta E_P/2)}
         {3\sinh(\beta E_K/2)}
    f(P)
  \right],
  \label{eq:stanford-1d}
\end{equation}
where \(m_1(P,K)\) is the angle-averaged kernel obtained from
Eq.~\eqref{eq:stanford-mkp} and Eq.~\eqref{eq:angular-measure}. In the scalar
benchmark the explicit expression is
\begin{equation}
  m_1(P,K)
  =
  \frac{3\lambda^2}{(2\pi)^3\beta}
  \frac{K}{P E_P E_K}
  \int_{|K-P|}^{K+P}\dd y\,
  \mathcal B(P,K,y),
  \label{eq:m1}
\end{equation}
with
\begin{align}
  \mathcal B(P,K,y)
  &=
  \frac{
    \log\!\left[\sinh x^{+}_{+}/\sinh x^{+}_{-}\right]
  }
  {\sinh(\beta E_+/2)}
  +
  \frac{
    \log\!\left[(1-e^{-2x^{-}_{+}})/(1-e^{2x^{-}_{-}})\right]
  }
  {\sinh(\beta E_-/2)},                                    \\
  E_\pm &= |E_K\pm E_P|,                                    \\
  x^{+}_{\pm}
  &=
  \frac{\beta}{4}
  \left[
    E_+
    \pm y\sqrt{1+\frac{4m^2}{y^2-E_+^2}}
  \right],                                                   \\
  x^{-}_{\pm}
  &=
  \frac{\beta}{4}
  \left[
    E_-
    \pm y\sqrt{1+\frac{4m^2}{y^2-E_-^2}}
  \right].
  \label{eq:xpm-1d}
\end{align}
The apparent \(E_-=0\) singularity must be handled by an analytic limit, not by
an arbitrary regulator.

Map the half-line to the unit interval:
\begin{equation}
  P(u)=P_0\frac{u}{1-u},
  \qquad 0<u<1.
  \label{eq:half-line-map}
\end{equation}
Use \(P_0\simeq 3m\) in the small-mass benchmark and \(P_0\simeq m\) away from
that regime. Define
\begin{equation}
  f_2(u)=(1-u)^{-1}P(u)f(P(u)),
  \label{eq:f2}
\end{equation}
\begin{equation}
  m_2(u,v)
  =
  P_0
  \frac{(1-u)(1-v)}{P(u)K(v)}
  m_1(P(u),K(v)),
  \label{eq:m2}
\end{equation}
and
\begin{equation}
  D(u)
  =
  \frac{P(u)}{1-u}
  \sinh\frac{\beta E_{P(u)}}{2}.
  \label{eq:D}
\end{equation}
Then
\begin{equation}
  \lambda_L f_2(u)
  =
  \int_0^1\dd v\,m_2(u,v)
  \left[
    f_2(v)-\frac{D(v)}{3D(u)}f_2(u)
  \right].
  \label{eq:symmetric-equation}
\end{equation}
After quadrature with nodes \(u_i\) and weights \(w_i\), assemble
\begin{equation}
  M_{ij}
  =
  w_j m_2(u_i,u_j)
  -
  \delta_{ij}
  \sum_{\ell}
  w_\ell m_2(u_i,u_\ell)
  \frac{D(u_\ell)}{3D(u_i)}.
  \label{eq:discrete-matrix}
\end{equation}
The largest positive eigenvalue of \(M\) is the numerical estimate of
\(\lambda_L\). The code should symmetrize only after applying the mathematically
correct similarity/weight transform; otherwise it should explicitly report the
antisymmetric residual.

\section{Leading-Behavior Extraction}

The extractor should keep the symbolic prefactor of each kernel contribution.
For a weak-coupling family of kernels,
\begin{equation}
  M
  =
  \sum_\gamma
  c_\gamma
  \beta^{-1}
  \prod_\alpha
  \bigl(g_\alpha\beta^{\Delta_\alpha}\bigr)^{n_{\gamma\alpha}}
  (\beta m_{\eff})^{-r_\gamma}
  \widehat M_\gamma(\beta m_{\eff},\hbox{dimensionless ratios}),
  \label{eq:scaling-kernel}
\end{equation}
the leading term is selected by the requested asymptotic ordering of
\(g_\alpha\), \(1/N\), and \(\beta m_{\eff}\). Numerically, this is checked by
log--log fits:
\begin{equation}
  \nu_g =
  \frac{\partial\log\lambda_L}{\partial\log g},
  \qquad
  \nu_m =
  \frac{\partial\log\lambda_L}{\partial\log m_{\eff}},
  \qquad
  \nu_\beta =
  \frac{\partial\log\lambda_L}{\partial\log \beta}.
  \label{eq:exponent-fits}
\end{equation}

For the Stanford benchmark,
\begin{equation}
  \lambda_L
  =
  c_m\,\frac{\lambda^2}{\beta^2 m}
  +\cdots,
  \qquad
  c_m\simeq 0.025,
  \qquad m\beta\ll 1,\ m\neq 0.
  \label{eq:stanford-leading-massive}
\end{equation}
For the massless theory,
\begin{equation}
  m_{\eff}=m_{\thm}
  =
  \sqrt{\frac{2\lambda}{3\beta^2}},
  \qquad
  \lambda_L
  =
  c_m\sqrt{\frac{3}{2}}\,
  \frac{\lambda^{3/2}}{\beta}
  \simeq
  0.031\,\frac{\lambda^{3/2}}{\beta}.
  \label{eq:stanford-leading-massless}
\end{equation}

\section{Implementation Algorithm}

\begin{enumerate}
  \item Parse the model into the data of Eq.~\eqref{eq:model}.
  \item Build \(S_{\SK}\) by Eq.~\eqref{eq:sk-action}; verify
  Eq.~\eqref{eq:unitarity-check} and the odd-\(a\) selection rule.
  \item Derive SK vertices from Eq.~\eqref{eq:sk-vertex}; store branch
  coefficients separately from couplings, signs, and index tensors.
  \item Derive \(G_{\R}\), \(\rho\), and \(\widetilde G\) from
  Eqs.~\eqref{eq:spectral-retarded}--\eqref{eq:half-wightman}.
  \item Compute required thermal masses and widths using
  Eqs.~\eqref{eq:thermal-mass-general} and~\eqref{eq:width}.
  \item Enumerate leading rung graphs and construct \(R_{AB}\) from
  Eq.~\eqref{eq:general-rung}.
  \item Put rails on shell and assemble \(M\) using
  Eqs.~\eqref{eq:gain-kernel}--\eqref{eq:M-general}.
  \item Reduce by symmetries: rotational invariance, species blocks, charge
  sectors, and large-\(N\) planar sectors.
  \item Discretize the resulting integral equation. For the scalar benchmark,
  use Eqs.~\eqref{eq:half-line-map}--\eqref{eq:discrete-matrix}.
  \item Extract \(\lambda_L\), eigenvectors, scaling exponents, and benchmark
  residuals.
\end{enumerate}

\section{Required Checks}

The implementation should fail closed if any of these checks fails:
\begin{itemize}
  \item \(S_{\SK}[\phi_a=0]=0\).
  \item No interaction vertex has an even number of \(a\)-legs.
  \item KMS relation: \(\widetilde G(k)=\rho(k)/(2\sinh(\beta k^0/2))\) for
  bosons.
  \item Rung support and theta-function regions agree with the on-shell
  kinematics.
  \item Widths satisfy \(\Gamma_a(\mathbf p)\ge 0\).
  \item The discretized operator reproduces the Stanford coefficients
  \(0.025\) and \(0.031\) within stated numerical tolerance.
  \item Leading exponents from symbolic prefactors agree with numerical
  log--log fits.
\end{itemize}

\begin{thebibliography}{9}
\bibitem{Stanford:2015owe}
Douglas Stanford,
\emph{Many-body chaos at weak coupling},
arXiv:1512.07687,
\href{https://arxiv.org/abs/1512.07687}{https://arxiv.org/abs/1512.07687}.
\end{thebibliography}

\end{document}
